\documentclass[11pt, a4paper]{article}

% ============================================================
% PACKAGES
% ============================================================
\usepackage[utf8]{inputenc}
\usepackage[T1]{fontenc}
\usepackage[french]{babel}
\usepackage{amsmath, amssymb}
\usepackage{geometry}
\usepackage{graphicx}
\usepackage{listings}
\usepackage{xcolor}
\usepackage{hyperref}
\usepackage{booktabs}
\usepackage{enumitem}
\usepackage{tcolorbox}
\usepackage{fancyhdr}
\usepackage{titlesec}

\geometry{margin=2.5cm}

% ============================================================
% STYLE
% ============================================================
\hypersetup{
    colorlinks=true,
    linkcolor=blue!60!black,
    citecolor=blue!60!black,
    urlcolor=blue!60!black
}

\definecolor{codegreen}{rgb}{0.15,0.55,0.15}
\definecolor{codegray}{rgb}{0.5,0.5,0.5}
\definecolor{codepurple}{rgb}{0.58,0,0.82}
\definecolor{backcolour}{rgb}{0.97,0.97,0.97}
\definecolor{accentblue}{rgb}{0.1,0.3,0.6}

\lstdefinestyle{pythonstyle}{
    backgroundcolor=\color{backcolour},
    commentstyle=\color{codegreen}\itshape,
    keywordstyle=\color{accentblue}\bfseries,
    numberstyle=\tiny\color{codegray},
    stringstyle=\color{codepurple},
    basicstyle=\ttfamily\small,
    breaklines=true,
    captionpos=b,
    keepspaces=true,
    numbers=left,
    numbersep=8pt,
    showspaces=false,
    showstringspaces=false,
    tabsize=4,
    language=Python,
    frame=single,
    rulecolor=\color{gray!30},
    morekeywords={np, pd, njit, NDArray}
}
\lstset{style=pythonstyle}

\tcbuselibrary{skins, breakable}
\newtcolorbox{keybox}[1][]{
    colback=blue!5!white,
    colframe=accentblue,
    fonttitle=\bfseries,
    title=#1,
    breakable
}
\newtcolorbox{warnbox}[1][]{
    colback=red!5!white,
    colframe=red!60!black,
    fonttitle=\bfseries,
    title=#1,
    breakable
}
\newtcolorbox{tipbox}[1][]{
    colback=green!5!white,
    colframe=green!50!black,
    fonttitle=\bfseries,
    title=#1,
    breakable
}

\pagestyle{fancy}
\fancyhf{}
\fancyhead[L]{\small\textit{Optimisation Python Num\'erique}}
\fancyhead[R]{\small\thepage}
\renewcommand{\headrulewidth}{0.4pt}

% ============================================================
% DOCUMENT
% ============================================================
\begin{document}

\begin{center}
    {\LARGE\bfseries Optimisation du Code Python Num\'erique}\\[0.4cm]
    {\large Retour d'exp\'erience sur un pipeline de Finance Quantitative}\\[0.8cm]
    {\normalsize Mars 2026}\\[0.3cm]
    \rule{\textwidth}{0.4pt}
\end{center}

\tableofcontents
\newpage

% ============================================================
\section{Introduction : Pourquoi optimiser ?}
% ============================================================

En Python scientifique, la performance repose sur un principe fondamental :

\begin{keybox}[R\`egle d'or]
\textbf{Minimiser le temps pass\'e dans l'interpr\'eteur Python.}\\[0.2cm]
Chaque boucle \texttt{for} en Python pur co\^ute $\sim$50--100\,ns par it\'eration.
La m\^eme op\'eration en NumPy co\^ute $\sim$1--5\,ns (vectoris\'ee sur le CPU).
Le facteur est de \textbf{10x \`a 100x}.
\end{keybox}

Ce document pr\'esente les techniques d'optimisation appliqu\'ees \`a un pipeline r\'eel
de traitement de donn\'ees tick haute fr\'equence (fichiers Parquet de $\sim$10\,GB),
en conservant la lisibilit\'e et la maintenabilit\'e du code.


% ============================================================
\section{Niveau 1 --- Vectorisation NumPy}
% ============================================================

\subsection{Principe}

La vectorisation consiste \`a remplacer les boucles Python par des op\'erations
sur des tableaux entiers. NumPy d\'el\`egue le calcul \`a des routines C/Fortran
optimis\'ees (BLAS, LAPACK, SIMD).

\begin{equation}
    \text{VWAP} = \frac{\sum_{i=1}^{N} p_i \cdot v_i}{\sum_{i=1}^{N} v_i}
\end{equation}

\subsection{Exemple : Tick Bars avec \texttt{reshape} + \texttt{einsum}}

Quand le nombre de ticks par bar est fixe ($T$), on peut reformater
les donn\'ees en matrice $(N_{\text{bars}} \times T)$ et agr\'eger sur l'axe 1 :

\begin{lstlisting}
# Reshape en matrice (n_bars x threshold)
price_matrix = arr_price.reshape(n_bars, threshold)
size_matrix  = arr_size.reshape(n_bars, threshold)

# Aggregations vectorisees
vol_sum = size_matrix.sum(axis=1)
dollar  = np.einsum("ij,ij->i", price_matrix, size_matrix)

# OHLCV en une seule passe
bars = pd.DataFrame({
    'open':  price_matrix[:, 0],
    'high':  price_matrix.max(axis=1),
    'low':   price_matrix.min(axis=1),
    'close': price_matrix[:, -1],
    'vwap':  dollar / vol_sum,
    'volume': vol_sum,
})
\end{lstlisting}

\begin{keybox}[\texttt{np.einsum} --- Notation d'Einstein]
\texttt{np.einsum("ij,ij->i", A, B)} calcule $\sum_j A_{ij} \cdot B_{ij}$
pour chaque ligne $i$. C'est \'equivalent \`a \texttt{(A * B).sum(axis=1)} mais
sans cr\'eer la matrice interm\'ediaire \texttt{A * B} en m\'emoire.
\end{keybox}


\subsection{Exemple : \texttt{np.reduceat} pour groupes de taille variable}

Quand les bars ont des tailles variables (Imbalance Bars, Runs Bars),
on ne peut pas utiliser \texttt{reshape}. La solution :
\texttt{np.ufunc.reduceat}.

\begin{lstlisting}
# bar_indices = [end_0, end_1, ..., end_k]
# starts      = [0, end_0+1, end_1+1, ...]

starts = np.empty(len(bar_indices), dtype=np.intp)
starts[0] = 0
starts[1:] = bar_indices[:-1] + 1

# Aggregations par groupe, SANS groupby pandas
high_   = np.maximum.reduceat(prices, starts)
low_    = np.minimum.reduceat(prices, starts)
volume  = np.add.reduceat(sizes, starts)
dollar  = np.add.reduceat(prices * sizes, starts)
\end{lstlisting}

\begin{keybox}[\texttt{np.add.reduceat} --- Comment \c{c}a marche]
Soit un tableau $x$ et des indices de d\'epart $[s_0, s_1, \ldots, s_k]$.\\
\texttt{np.add.reduceat(x, starts)} retourne :
\[
    \texttt{out}[i] = \sum_{j=s_i}^{s_{i+1}-1} x[j]
\]
pour chaque segment. Fonctionne avec \texttt{add}, \texttt{maximum}, \texttt{minimum},
et tout \texttt{ufunc} ayant une op\'eration d'identit\'e.
\end{keybox}

\subsubsection{Comparaison de performance : \texttt{groupby} vs \texttt{reduceat}}

\begin{center}
\begin{tabular}{lccc}
\toprule
\textbf{M\'ethode} & \textbf{10M ticks} & \textbf{100M ticks} & \textbf{Allocation m\'emoire} \\
\midrule
\texttt{pd.groupby().agg()} & $\sim$2\,s & $\sim$25\,s & $\mathcal{O}(N)$ copies \\
\texttt{np.reduceat} & $\sim$0.3\,s & $\sim$3\,s & $\mathcal{O}(K)$ (nb bars) \\
\bottomrule
\end{tabular}
\end{center}

Le gain provient de deux facteurs :
\begin{enumerate}
    \item \textbf{Z\'ero overhead Python} : pas de dispatch pandas, pas de MultiIndex.
    \item \textbf{Z\'ero copie} : \texttt{reduceat} op\`ere sur des vues du tableau original.
\end{enumerate}


% ============================================================
\section{Niveau 2 --- R\'eduction des allocations m\'emoire}
% ============================================================

\subsection{Le probl\`eme : chaque copie co\^ute cher}

En Python, chaque \texttt{.copy()}, chaque cr\'eation de \texttt{pd.Series}
interm\'ediaire, et chaque appel \`a \texttt{.replace()} alloue de la m\'emoire
sur le heap. Pour un tableau de 100M \'el\'ements float64 :

\begin{equation}
    100 \times 10^6 \times 8\,\text{octets} = 800\,\text{MB par copie}
\end{equation}

\subsection{Exemple : \texttt{tick\_rule\_creator}}

\subsubsection{Version originale (2 copies interm\'ediaires)}
\begin{lstlisting}
# Version originale : 2 copies intermediaires
diff_sign = np.sign(price_series.diff())        # copie 1
diff_sign = diff_sign.replace(0, np.nan).ffill() # copie 2
diff_sign.iloc[0] = np.nan                       # mutation
\end{lstlisting}

\subsubsection{Version optimis\'ee (1 allocation contr\^ol\'ee)}
\begin{lstlisting}
# Version optimisee : 1 seule allocation
prices = np.asarray(price_series, dtype=np.float64)
diff = np.empty(len(prices), dtype=np.float64)
diff[0] = np.nan
diff[1:] = np.sign(np.diff(prices))

# In-place : pas de copie supplementaire
diff[diff == 0.0] = np.nan
return pd.Series(diff).ffill().to_numpy(dtype=np.float64)
\end{lstlisting}

\begin{tipbox}[Astuce : \texttt{np.empty} vs \texttt{np.zeros}]
\texttt{np.empty(n)} alloue la m\'emoire \textbf{sans l'initialiser} (plus rapide).
\`A utiliser quand on \'ecrit imm\'ediatement dans chaque case.
\texttt{np.zeros(n)} initialise \`a z\'ero --- n\'ecessaire si certaines cases
ne sont pas explicitement remplies.
\end{tipbox}


\subsection{Fusionner les \texttt{groupby} redondants}

Chaque \texttt{groupby} pandas reconstruit un dictionnaire de groupes.
Si on en fait deux sur le m\^eme \texttt{bar\_id}, on double le co\^ut :

\begin{lstlisting}
# AVANT : 2 groupby sur le meme bar_id
bars = data.groupby('bar_id').agg({...})      # groupby 1
vwap = data.groupby('bar_id')['dollar'].sum() # groupby 2

# APRES : 1 seul groupby
data["dollar_value"] = data['price'] * data['size']
bars = data.groupby('bar_id').agg({
    ...,
    'dollar_value': 'sum'   # inclus dans le meme appel
})
\end{lstlisting}

Gain : $\sim$\textbf{2x} sur la fonction concern\'ee.


% ============================================================
\section{Niveau 3 --- Numba pour les boucles path-dependent}
% ============================================================

\subsection{Quand utiliser Numba ?}

Certains algorithmes sont intrins\`equement s\'equentiels : chaque it\'eration
d\'epend du r\'esultat de la pr\'ec\'edente. C'est le cas des \textbf{EWMA r\'ecursives}
utilis\'ees dans les Imbalance/Runs Bars :

\begin{equation}
    \hat{E}_T^{(n)} = (1 - \alpha)\,\hat{E}_T^{(n-1)} + \alpha\,T_n
\end{equation}

Ces formules \textbf{ne peuvent pas \^etre vectoris\'ees} car $\hat{E}_T^{(n)}$
d\'epend de $\hat{E}_T^{(n-1)}$. NumPy ne peut pas aider ici.

\subsection{Numba JIT : le meilleur des deux mondes}

\begin{lstlisting}
from numba import njit

@njit(cache=True, fastmath=True)
def get_imbalance_bars_numba(tick_signs, roll_indices,
                              initial_T, initial_E_b, ...):
    T = len(tick_signs)
    bar_indices = np.zeros(T, dtype=np.int32)
    bar_count = 0
    theta = 0.0
    E_T = initial_T
    E_b = initial_E_b

    for i in range(T):
        theta += tick_signs[i]
        threshold = E_T * np.abs(E_b)

        if np.abs(theta) >= threshold:
            bar_indices[bar_count] = i
            bar_count += 1
            # EWMA update (path-dependent)
            current_T = i - last_idx
            E_T = (1 - alpha) * E_T + alpha * current_T
            ...
    return bar_indices[:bar_count]
\end{lstlisting}

\begin{keybox}[Options \texttt{@njit}]
\begin{description}[leftmargin=2cm]
    \item[\texttt{cache=True}] Sauvegarde le bytecode compil\'e sur disque.
        \'Evite la recompilation \`a chaque import ($\sim$1--3\,s \'economis\'es).
    \item[\texttt{fastmath=True}] Autorise les r\'eorganisations flottantes
        (associativit\'e, FMA). Gain de $\sim$10--20\% sur les boucles arithm\'etiques.
        \textbf{Attention} : peut changer les r\'esultats au niveau de l'ULP
        (Unit in the Last Place).
\end{description}
\end{keybox}

\subsection{Arbre de d\'ecision : quelle technique choisir ?}

\begin{center}
\begin{tabular}{p{5cm}p{4cm}p{4cm}}
\toprule
\textbf{Situation} & \textbf{Technique} & \textbf{Gain typique} \\
\midrule
Op\'eration \'el\'ement par \'el\'ement & NumPy vectoris\'e & 10--100x \\
Agr\'egation par groupes (taille fixe) & \texttt{reshape} + \texttt{axis} & 5--20x \\
Agr\'egation par groupes (taille variable) & \texttt{np.reduceat} & 5--10x \\
Boucle avec d\'ependance s\'equentielle & Numba \texttt{@njit} & 50--200x \\
Parall\'elisme multi-fichiers & \texttt{ProcessPoolExecutor} & $\sim N_{\text{cores}}$x \\
\bottomrule
\end{tabular}
\end{center}


% ============================================================
\section{Niveau 4 --- Pi\`eges de pandas 2.x / PyArrow}
% ============================================================

\subsection{Copy-on-Write (CoW)}

Depuis pandas 2.0, le mode \textbf{Copy-on-Write} est activ\'e par d\'efaut.
Cons\'equence : \texttt{.values} retourne souvent un array \textbf{read-only}.

\begin{warnbox}[Erreur classique : \texttt{assignment destination is read-only}]
\begin{lstlisting}
result = df['price'].values   # read-only avec CoW
result[0] = np.nan            # ERREUR !
\end{lstlisting}
\end{warnbox}

\subsubsection{Solutions}

\begin{center}
\begin{tabular}{lp{7cm}}
\toprule
\textbf{M\'ethode} & \textbf{Comportement} \\
\midrule
\texttt{.values} & Retourne une vue (potentiellement read-only avec CoW) \\
\texttt{.to\_numpy()} & Retourne un array NumPy (writable dans la plupart des cas) \\
\texttt{.to\_numpy(dtype=..., copy=True)} & Force une copie writable \\
\texttt{np.asarray(series, dtype=...)} & Conversion + copie si le dtype diff\`ere \\
\bottomrule
\end{tabular}
\end{center}

\begin{tipbox}[R\`egle simple]
\begin{itemize}
    \item Lecture seule $\rightarrow$ \texttt{.values} (z\'ero copie, plus rapide)
    \item \'Ecriture n\'ecessaire $\rightarrow$ \texttt{.to\_numpy(dtype=np.float64)}
\end{itemize}
\end{tipbox}

\subsection{Backend PyArrow et arrays read-only}

Quand un DataFrame est charg\'e depuis Parquet avec \texttt{engine="pyarrow"},
les colonnes sont stock\'ees en m\'emoire PyArrow (z\'ero copie depuis le fichier).
Ces buffers sont \textbf{immutables par d\'esign}.

\begin{lstlisting}
df = pd.read_parquet("data.parquet", engine="pyarrow")
arr = df['price'].values  # ArrowExtensionArray -> read-only

# Solution : forcer la conversion en NumPy writable
arr = np.asarray(df['price'], dtype=np.float64)  # copie writable
\end{lstlisting}


% ============================================================
\section{Niveau 5 --- Factorisation et DRY}
% ============================================================

\subsection{Principe DRY : Don't Repeat Yourself}

Quand $N$ fonctions partagent le m\^eme bloc de pr\'eparation, chaque
modification doit \^etre r\'epliqu\'ee $N$ fois. C'est une source de bugs
\textbf{et} de perte de performance (on oublie d'optimiser une copie).

\subsection{Exemple : \texttt{\_prepare\_data}}

8 fonctions (3 imbalance + 3 runs + 2 variantes) partageaient :

\begin{lstlisting}
# Code duplique dans 8 fonctions
data = ticks_dataset[["ts_event", "price", "size"]].copy()
data["tick_rule"] = tick_rule_creator(data["price"])
data.dropna(subset=["tick_rule"], inplace=True)
\end{lstlisting}

Factoris\'e en un seul helper :

\begin{lstlisting}
def _prepare_data(ticks_dataset: pd.DataFrame) -> pd.DataFrame:
    """Preparation commune a tous les bar creators."""
    data = ticks_dataset[["ts_event", "price", "size"]].copy()
    data["tick_rule"] = tick_rule_creator(data["price"])
    data.dropna(subset=["tick_rule"], inplace=True)
    return data
\end{lstlisting}

B\'en\'efices :
\begin{enumerate}
    \item \textbf{Une seule optimisation} profite \`a toutes les fonctions.
    \item \textbf{Un seul point de correction} en cas de bug (ex: CoW read-only).
    \item \textbf{Moins de code} = moins de surface d'attaque pour les erreurs.
\end{enumerate}


% ============================================================
\section{Niveau 6 --- Parall\'elisme multi-processus}
% ============================================================

\subsection{Quand parall\'eliser ?}

Le parall\'elisme via \texttt{ProcessPoolExecutor} est pertinent quand :
\begin{enumerate}
    \item Les t\^aches sont \textbf{ind\'ependantes} (pas de donn\'ees partag\'ees).
    \item Le temps de calcul par t\^ache $\gg$ le co\^ut de s\'erialisation pickle.
    \item Les donn\'ees en entr\'ee sont \textbf{grosses} ($> 100$\,MB par asset).
\end{enumerate}

\subsection{Co\^ut cach\'e : la s\'erialisation}

Chaque worker re\c{c}oit ses arguments \textbf{s\'erialis\'es via pickle}.
Pour un DataFrame de 1\,GB :

\begin{equation}
    t_{\text{pickle}} \approx 2\text{--}5\,\text{s (s\'erialisation + d\'es\'erialisation)}
\end{equation}

\begin{tipbox}[Astuce : passer le chemin, pas les donn\'ees]
Au lieu de passer le DataFrame au worker, on passe le \textbf{chemin du fichier}.
Chaque worker charge lui-m\^eme le Parquet (z\'ero s\'erialisation).

\begin{lstlisting}
# BON : chaque worker charge son propre fichier
def process_one_asset(asset, datadir, month, config):
    df = data_loader(datadir, asset, month)  # I/O local
    ...
\end{lstlisting}
\end{tipbox}

\subsection{Nombre optimal de workers}

\begin{equation}
    N_{\text{workers}} = \min\bigl(N_{\text{assets}},\; N_{\text{CPU}}\bigr)
\end{equation}

Au-del\`a de $N_{\text{CPU}}$, les workers se disputent le CPU (context switching)
et la bande passante m\'emoire (cache thrashing).


% ============================================================
\section{R\'esum\'e : checklist d'optimisation}
% ============================================================

\begin{keybox}[Checklist \`a suivre dans l'ordre]
\begin{enumerate}[leftmargin=1.5cm]
    \item[\textbf{1.}] \textbf{Profiler d'abord} --- Ne jamais optimiser \`a l'aveugle.
        Utiliser \texttt{cProfile}, \texttt{line\_profiler}, ou \texttt{\%timeit}.
    \item[\textbf{2.}] \textbf{Vectoriser} --- Remplacer les boucles par NumPy
        (\texttt{reduceat}, \texttt{einsum}, broadcasting).
    \item[\textbf{3.}] \textbf{R\'eduire les copies} --- \'Eliminer les allocations
        interm\'ediaires (\texttt{.replace()}, \texttt{.copy()} inutiles).
    \item[\textbf{4.}] \textbf{Numba pour le s\'equentiel} --- \texttt{@njit} uniquement
        pour les boucles path-dependent (EWMA, theta r\'ecursif).
    \item[\textbf{5.}] \textbf{Factoriser} --- Un seul point d'optimisation via DRY.
    \item[\textbf{6.}] \textbf{Parall\'eliser en dernier} --- Seulement si les t\^aches
        sont ind\'ependantes et les donn\'ees volumineuses.
\end{enumerate}
\end{keybox}

\vfill
\begin{center}
\rule{0.5\textwidth}{0.4pt}\\[0.3cm]
{\small\textit{Document g\'en\'er\'e dans le cadre du projet ML Finance --- Mars 2026}}
\end{center}

\end{document}
